\documentclass[11pt]{article}
\usepackage[a4paper,margin=25mm]{geometry}
\usepackage{amsmath,amssymb,amsthm}
\usepackage[hidelinks]{hyperref}
\emergencystretch=2em
\newtheorem{theorem}{Theorem}
\title{Systematic resampling can increase variance\\
\large Disproof of Conjecture 00000008263}
\author{GitHub: gaochengzhecpu}\date{}
\begin{document}\maketitle
\begin{abstract}
Four positively weighted particles give a systematic resampling estimator
with variance $1/16$, while multinomial resampling has variance $3/64$.
Thus systematic resampling can be worse, and cannot universally halve
multinomial variance, whether halving means equality or an upper bound.
The exact inverse-CDF algorithms and finite probability calculations
are verified in Lean.
\end{abstract}

\section*{The assertion being tested}
The source claims a variance-halving law for systematic versus multinomial
resampling. We test the usual conditional resampling variance of the
equally weighted empirical estimate of a fixed particle observable,
conditional on the input particles and normalized weights. A universal
variance-halving law would in particular require
\[
 \operatorname{Var}(\widehat\mu_{\rm sys}(f))
 \leq\tfrac12\operatorname{Var}(\widehat\mu_{\rm mult}(f))
\]
for every such input and observable. We give a strict counterexample even
to the weaker assertion that systematic variance never exceeds multinomial
variance. This refutes the variance conjunct; it is unnecessary to settle
the source's effective-sample-size or other error-constant assertions.

This phenomenon is known: systematic resampling does not enjoy a universal
variance improvement over multinomial resampling; see Douc, Capp\'e and
Moulines, \emph{Comparison of Resampling Schemes for Particle Filtering}
(2005), \href{https://arxiv.org/abs/cs/0507025}{arXiv:cs/0507025}.
The present four-particle calculation is supplied as an explicit,
independently checkable instance, not as a new general phenomenon.

\section*{Particles, ordering, and the actual algorithms}
Keep the given particle ordering and choose
\[
 (w_0,w_1,w_2,w_3)=(1/8,3/8,1/8,3/8),\qquad
 (f_0,f_1,f_2,f_3)=(1,0,1,0).
\]
All weights are strictly positive and their sum is one. The target mean
is $\mu(f)=\sum_i w_i f_i=1/4$.
The inverse CDF sends a position $x\in[0,1)$ to the unique particle
whose half-open cumulative-weight interval contains $x$:
\[
 L(x)=\begin{cases}
 0,&0\leq x<1/8,\\
 1,&1/8\leq x<1/2,\\
 2,&1/2\leq x<5/8,\\
 3,&5/8\leq x<1.
 \end{cases}
\]
Systematic resampling chooses one uniform offset $U\in[0,1/4)$ and
selects the labels $L(U+k/4)$, $k=0,1,2,3$. Its estimator is
\[
 \widehat\mu_{\rm sys}(f)=\tfrac14\sum_{k=0}^3 f_{L(U+k/4)}.
\]
Multinomial resampling chooses four independent labels $J_k$, each with
$\mathbb P(J_k=i)=w_i$, and uses
$\widehat\mu_{\rm mult}(f)=\tfrac14\sum_{k=0}^3 f_{J_k}$.
The particles may be placed at four distinct state-space points with the
displayed observable values, so no coincidence or particle identification
is needed.

\section*{Exact systematic variance}
For every offset, including the interval boundaries under the stated
half-open convention, the selected label vector is
\[
 (L(U),L(U+1/4),L(U+1/2),L(U+3/4))
 =\begin{cases}
 (0,1,2,3),&0\leq U<1/8,\\
 (1,1,3,3),&1/8\leq U<1/4.
 \end{cases}
\]
Each phase interval has probability $1/2$, since it has length $1/8$
inside an interval of length $1/4$. Hence
\[
 \widehat\mu_{\rm sys}(f)=
 \begin{cases}1/2,&\text{with probability }1/2,\\
 0,&\text{with probability }1/2.
 \end{cases}
\]
It is unbiased and has variance
\[
 \tfrac12(1/2-1/4)^2+\tfrac12(0-1/4)^2=\frac1{16}.
\]

\section*{Exact multinomial variance and contradiction}
Each $f_{J_k}$ is a Bernoulli variable with success probability
$w_0+w_2=1/4$. Independence gives
\[
 \operatorname{Var}(\widehat\mu_{\rm mult}(f))
 =\frac1{4^2}\sum_{k=0}^3\frac14\left(1-\frac14\right)
 =\frac3{64}.
\]
Both estimators therefore have the same correct expectation $1/4$, but
\[
 \operatorname{Var}(\widehat\mu_{\rm sys}(f))=\frac4{64}
 >\frac3{64}=\operatorname{Var}(\widehat\mu_{\rm mult}(f)).
\]
In particular systematic variance is neither equal to nor bounded by
one half of multinomial variance. The same strict comparison holds for
the unnormalized sum estimator, since both variances are multiplied by
$16$. An extra assumption imposing a special particle ordering or an
additional random permutation would describe a different claim. No such
assumption appears in the conjecture; systematic resampling in the stated
ordering is the standard algorithm.

\section*{Formal verification and scope}
Lean defines the actual four-particle inverse-CDF algorithm from its
cumulative weights and uses the actual systematic positions $U+k/4$.
For an arbitrary ordered field, it proves the selected-label formula
above for \emph{every} offset in $[0,1/4)$. Thus it applies in particular
to real offsets, and does not merely evaluate two chosen offsets.

A finite probability law is defined by exact rational masses, with
proofs of nonnegativity and total mass one. Expectation is the
mass-weighted finite sum and variance is the expectation of the squared
deviation from the expectation. The systematic law is uniform on the two
phase strata; the estimator is still generated by inverse CDF at a
representative of each stratum. The selected-label theorem justifies this
representation for every observable. The multinomial law is the genuine
product law on all $4^4=256$ ordered label tuples, with mass equal to the
product of their four input weights.

Lean computes the expectations and centered variances from these
algorithms and laws, proves the strict worse-than-multinomial comparison,
and negates the necessary universal halving assertion over all rational
observables on this fixed legal weighted population. The final theorem is
\begin{center}\texttt{Conjecture8263.conjecture8263\_false}.\end{center}
The counterexample observable is rational, so restriction to rational
observables cannot hide the failure of a claim over real observables.

The elementary measure identification of a continuous uniform offset
with the two equal-length phase strata is supplied on paper. Lean
formalizes all offsets' algorithmic outputs and the resulting finite
probability calculations, not a general measure-theoretic implementation
of uniform real random variables or particle-filter evolution. The
finite sums are computations of the actual algorithmic random variables,
not an imposed table of variance values.

\section*{Source and reproduction}
The unmodified bilingual source is supplied as \texttt{SOURCE.md},
with frozen-source provenance and SHA-256 in the verification folder.
The coordinating workflow found no prior submission for this conjecture.
Lean is pinned to 4.19.0 and Mathlib commit
\begin{center}\texttt{c44e0c8ee63ca166450922a373c7409c5d26b00b}.\end{center}
Run in the supplied Lean folder:
\begin{center}\texttt{lake env lean -DwarningAsError=true Main.lean}.\end{center}
This draft was developed by a delegated AI agent with a self-review;
the coordinating agent reviews it separately. No external independent
review is claimed.
\end{document}
